\documentclass[11pt,a4paper]{article}

\usepackage[french]{babel}
\usepackage[utf8]{inputenc}
\usepackage[T1]{fontenc}
\usepackage[margin=2.2cm]{geometry}
\usepackage{xcolor}
\usepackage{titlesec}
\usepackage{enumitem}
\usepackage{tabularx}
\usepackage{booktabs}
\usepackage{colortbl}
\usepackage{fancyhdr}
\usepackage{hyperref}
\usepackage{parskip}

\definecolor{titlecolor}{HTML}{1A1A2E}
\definecolor{h1color}{HTML}{16213E}
\definecolor{h2color}{HTML}{0F3460}
\definecolor{notebg}{HTML}{F0F4F8}
\definecolor{tablehead}{HTML}{0F3460}
\definecolor{rowalt}{HTML}{F5F7FA}
\definecolor{graytext}{HTML}{555555}

\titleformat{\section}{\Large\bfseries\color{h1color}}{\thesection}{1em}{}
\titleformat{\subsection}{\large\bfseries\color{h2color}}{\thesubsection}{1em}{}
\titlespacing*{\section}{0pt}{18pt}{8pt}
\titlespacing*{\subsection}{0pt}{12pt}{6pt}

\pagestyle{fancy}
\fancyhf{}
\fancyfoot[C]{\thepage}
\renewcommand{\headrulewidth}{0pt}

\newenvironment{note}
  {\begin{center}\begin{minipage}{0.96\textwidth}\colorbox{notebg}{\begin{minipage}{0.94\textwidth}\vspace{4pt}}
  }
  {\vspace{4pt}\end{minipage}\end{minipage}\end{center}}

\newcommand{\code}[1]{\texttt{#1}}

\title{
  \color{titlecolor}\textbf{\huge Prototype2Code}\\[4pt]
  \color{titlecolor}\Large De la maquette UI à l'Arbre de Mise en Page
}
\author{}
\date{}

\begin{document}

\begin{center}
  {\huge\bfseries\color{titlecolor} Prototype2Code}\\[6pt]
  {\Large\color{h1color} De la maquette UI à l'Arbre de Mise en Page (UI Layout Tree)}\\[10pt]
  {\color{graytext}\itshape Explication de la méthode décrite en section 3.2 et formalisée par l'Algorithme 1 :\\
  construction de la hiérarchie, puis segmentation Row/Column pour le modèle Flexbox.}
\end{center}

\vspace{4pt}
\noindent\rule{\textwidth}{0.6pt}
\vspace{10pt}

L'objectif du processus est de transformer une maquette (des rectangles positionnés en pixels, sans structure logique) en un \textbf{Arbre de Mise en Page (UI Layout Tree)} organisé en groupes imbriqués, façon « poupées russes ». Cette structure porte ensuite directement les instructions nécessaires à la génération de code HTML (imbrication de \code{<div>}) et CSS (Flexbox ou positionnement absolu). Le processus se déroule en deux grandes phases.

\section*{Phase 1 --- Initialisation de l'arbre et gestion de la hiérarchie}

Avant de construire l'arbre, le système exploite les résultats de l'étape de « linting » (correction) pour préparer des données propres.

\subsection*{1. Reconstruction des relations de contenance}
\begin{itemize}[leftmargin=1.2em, itemsep=4pt]
  \item Les groupes perceptuels détectés (cartes, barres de navigation, etc.) sont triés par taille ascendante, du plus petit au plus grand.
  \item Pour chaque couche, on calcule l'IoU (Intersection sur Union) afin de vérifier si elle appartient géométriquement à un groupe~; si oui, elle lui est rattachée.
  \item Les éléments fragmentés (ex. les morceaux d'une icône coupée par la détection) sont fusionnés en une seule entité avec de nouvelles coordonnées globales, du point haut-gauche au point bas-droit.
\end{itemize}

\subsection*{2. Construction de l'arbre initial (Algorithme 1, Phase 1)}
\begin{itemize}[leftmargin=1.2em, itemsep=4pt]
  \item Extraction de toutes les couches « feuilles » --- les éléments terminaux comme le texte ou les images.
  \item Ces éléments sont triés par taille \textbf{descendante}, du plus grand au plus petit.
  \item Chaque élément est inséré sous le nœud qui le contient physiquement et qui possède la plus petite taille possible --- garantissant une hiérarchie précise, du type « poupée russe ».
\end{itemize}

\subsection*{3. Gestion des chevauchements}
\begin{itemize}[leftmargin=1.2em, itemsep=4pt]
  \item Pour chaque nœud, le système vérifie si ses enfants se superposent.
  \item En cas de chevauchement, un nouveau nœud parent est créé avec l'attribut \code{need\_absolute=True}, indiquant que ces éléments devront utiliser un positionnement CSS absolu plutôt que le flux normal.
\end{itemize}



\section*{Phase 2 --- Segmentation Row/Column pour le modèle Flexbox}

Une fois la hiérarchie de base établie, il faut organiser les éléments de façon ordonnée pour permettre un design adaptatif (responsive).

\subsection*{1. Segmentation verticale (lignes)}
\begin{itemize}[leftmargin=1.2em, itemsep=4pt]
  \item L'algorithme cherche les coordonnées $Y$ où une ligne horizontale peut traverser toute la zone sans couper aucun élément.
  \item Cela divise l'espace en segments verticaux --- des « étages ».
\end{itemize}

\subsection*{2. Segmentation horizontale (colonnes)}
\begin{itemize}[leftmargin=1.2em, itemsep=4pt]
  \item À l'intérieur de chaque segment vertical, le système cherche des coordonnées $X$ pour tracer des lignes verticales sans couper d'éléments.
  \item Cela permet de regrouper les éléments côte à côte sur une même ligne.
\end{itemize}

\subsection*{3. Récursion et attribution Flexbox}
\begin{itemize}[leftmargin=1.2em, itemsep=4pt]
  \item Les deux étapes précédentes sont appliquées récursivement jusqu'à ce qu'il soit impossible de segmenter davantage.
  \item Si plusieurs éléments sont dans le même segment horizontal~: le parent reçoit \code{flex-direction: row}.
  \item S'ils sont alignés verticalement~: le parent reçoit \code{flex-direction: column}.
\end{itemize}

\section*{Synthèse des deux phases}

\begin{table}[h!]
\centering
\renewcommand{\arraystretch}{1.4}
\begin{tabularx}{\textwidth}{|>{\raggedright\arraybackslash}p{3.4cm}|>{\raggedright\arraybackslash}X|>{\raggedright\arraybackslash}X|}
\hline
\rowcolor{tablehead}
\textcolor{white}{\textbf{Phase}} & \textcolor{white}{\textbf{Objectif}} & \textcolor{white}{\textbf{Sortie produite}} \\
\hline
\rowcolor{white}
1. Initialisation de la hiérarchie &
Déterminer les relations de contenance et fusionner les éléments fragmentés &
Arbre brut « poupée russe » + nœuds \code{need\_absolute=True} \\
\hline
\rowcolor{rowalt}
2. Segmentation Flexbox &
Découper récursivement l'espace en lignes puis colonnes &
Arbre final avec \code{flex-direction: row / column} sur chaque parent \\
\hline
\end{tabularx}
\end{table}

\section*{Résultat final}

À l'issue de ce processus, on obtient un arbre structuré où chaque nœud porte des instructions explicites pour la génération de code~: imbrication HTML (des \code{<div>} logiquement groupés) et modèle de boîte CSS (Flexbox en \code{row}/\code{column}, ou positionnement absolu lorsque nécessaire). Cette organisation en « groupes imbriqués » rend le code final bien plus maintenable et lisible qu'un placement au pixel près.

\end{document}